% GENERATED FILE. Edit canonical lessons, then run npm run generate:books.
% canonical-source-hash: fnv1a64:50f48d90897d710e
% canonical-lessons: TA-C419-kulirutti, TA-C419-vennir, TA-C419-porttu, TA-C419-ilaipparu, TA-C419-viruntompal

\chapter{A Comfortable Room}
\label{ch:ta-a-comfortable-room}
\hlchaptermodality{419}

\begin{chapteropening}
\textbf{By the end of this chapter:} \emph{I can talk about an air conditioner, hot water and hospitality, and covering up and resting.}

Complete the last lesson of chapter 419: I can talk about an air conditioner, hot water and hospitality, and covering up and resting.
\end{chapteropening}

\section[\texorpdfstring{\ta{குளிரூட்டி} (kuḷirūṭṭi)}{குளிரூட்டி (kuḷirūṭṭi)}]{\ta{குளிரூட்டி} (kuḷirūṭṭi) — an air conditioner}

\label{lesson:TA-C419-kulirutti}



\begin{quote}
Before the new one: say the Tamil for bedding, then the Tamil for to spread out.
\end{quote}

\subsection*{You'll want to know: \ta{குளிரூட்டி}}
\textbf{\ta{குளிரூட்டி}} — \emph{kuḷirūṭṭi} — \textquotedblleft{}an air conditioner\textquotedblright{}.

\begin{quote}
\ta{அறையில்} \ta{குளிரூட்டி} \ta{இருந்ததால்} \ta{நன்றாகத்} \ta{தூங்கினோம்}.
\end{quote}

\begin{quote}
\emph{aṟaiyil kuḷirūṭṭi iruntatāl naṉṟākat tūṅkiṉōm.}
\end{quote}

\begin{quote}
Because there was an air conditioner in the room, we slept well.
\end{quote}

\ta{குளிர்}, cold, + \ta{ஊட்டி}, a feeder: a thing that feeds cold.

\begin{cousinweb}
\textbf{\ta{குளிரூட்டி}} (\emph{kuḷirūṭṭi}) is a modern coinage built from \ta{குளிர்}, \textquotedblleft{}cold\textquotedblright{} (DEDR 1834), and \ta{ஊட்டி}, \textquotedblleft{}a feeder\textquotedblright{}, from \ta{ஊட்டு}, \textquotedblleft{}to feed\textquotedblright{} (DEDR 600).
\end{cousinweb}

\begin{grammarlens}[title={What you've built}]
One more word for getting around.
\end{grammarlens}

\subsection*{Your turn}
\begin{itemize}
\raggedright
  \item \emph{Say it:} \emph{kuḷirūṭṭi}
  \item \emph{Say it:} \emph{kuḷirūṭṭi}, once more
  \item \emph{Say it:} \emph{aṟaiyil kuḷirūṭṭi iruntatāl naṉṟākat tūṅkiṉōm.} — Because there was an air conditioner in the room, we slept well.
  \item \emph{Recall:} say the Tamil for bedding, then the Tamil for to spread out, then say \emph{kuḷirūṭṭi} again
\end{itemize}

\begin{tcolorbox}[breakable,colback=teal!4,colframe=teal!35!black,title={Before you move on}]
What does \ta{குளிரூட்டி} mean? (\textquotedblleft{}An air conditioner\textquotedblright{}.) Say it once more.
\end{tcolorbox}
\section[\texorpdfstring{\ta{வெந்நீர்} (vennīr)}{வெந்நீர் (vennīr)}]{\ta{வெந்நீர்} (vennīr) — hot water}

\label{lesson:TA-C419-vennir}



\begin{quote}
Before the new one: say the Tamil for to spread out, then the Tamil for an air conditioner.
\end{quote}

\subsection*{You'll want to know: \ta{வெந்நீர்}}
\textbf{\ta{வெந்நீர்}} — \emph{vennīr} — \textquotedblleft{}hot water\textquotedblright{}.

\begin{quote}
\ta{குளிப்பதற்கு} \ta{வெந்நீர்} \ta{இல்லை}.
\end{quote}

\begin{quote}
\emph{kuḷippataṟku vennīr illai.}
\end{quote}

\begin{quote}
There is no hot water for bathing.
\end{quote}

\ta{வெம்}, hot, + \ta{நீர்}, water: the \ta{ம்} becomes \ta{ந்} before \ta{ந}.

\begin{cousinweb}
\textbf{\ta{வெந்நீர்}} (\emph{vennīr}) is native Dravidian: DEDR 5517 lists it, \textquotedblleft{}hot water\textquotedblright{}, from the root \ta{வே}, \textquotedblleft{}to burn, to be hot\textquotedblright{}, + \ta{நீர்}, \textquotedblleft{}water\textquotedblright{}.
\end{cousinweb}

\begin{grammarlens}[title={What you've built}]
One more word for getting around.
\end{grammarlens}

\subsection*{Your turn}
\begin{itemize}
\raggedright
  \item \emph{Say it:} \emph{vennīr}
  \item \emph{Say it:} \emph{vennīr}, once more
  \item \emph{Say it:} \emph{kuḷippataṟku vennīr illai.} — There is no hot water for bathing.
  \item \emph{Recall:} say the Tamil for to spread out, then the Tamil for an air conditioner, then say \emph{vennīr} again
\end{itemize}

\begin{tcolorbox}[breakable,colback=teal!4,colframe=teal!35!black,title={Before you move on}]
What does \ta{வெந்நீர்} mean? (\textquotedblleft{}Hot water\textquotedblright{}.) Say it once more.
\end{tcolorbox}
\section[\texorpdfstring{\ta{போர்த்து} (pōrttu)}{போர்த்து (pōrttu)}]{\ta{போர்த்து} (pōrttu) — to wrap up; to cover with a cloth}

\label{lesson:TA-C419-porttu}



\begin{quote}
Before the new one: say the Tamil for an air conditioner, then the Tamil for hot water.
\end{quote}

\subsection*{You'll want to know: \ta{போர்த்து}}
\textbf{\ta{போர்த்து}} — \emph{pōrttu} — \textquotedblleft{}to wrap up; to cover with a cloth\textquotedblright{}.

\begin{quote}
\ta{குளிரில்} \ta{நான்} \ta{குழந்தையைப்} \ta{போர்வையால்} \ta{போர்த்தினேன்}.
\end{quote}

\begin{quote}
\emph{kuḷiril nāṉ kuḻantaiyaip pōrvaiyāl pōrttiṉēṉ.}
\end{quote}

\begin{quote}
In the cold I wrapped the child in a blanket.
\end{quote}

Past (I …): \ta{போர்த்தினேன்}, \emph{pōrttiṉēṉ}.

\ta{போர்வை}, a blanket, comes from the same root.

\begin{cousinweb}
\textbf{\ta{போர்த்து}} (\emph{pōrttu}) is native Dravidian (DEDR 4590), from \ta{போர்}, \textquotedblleft{}to wrap, to cover\textquotedblright{}, the root of \ta{போர்வை}, \textquotedblleft{}a blanket\textquotedblright{}.
\end{cousinweb}

\begin{grammarlens}[title={What you've built}]
One more word for getting around.
\end{grammarlens}

\subsection*{Your turn}
\begin{itemize}
\raggedright
  \item \emph{Say it:} \emph{pōrttu}
  \item \emph{Say it:} \emph{pōrttu}, once more
  \item \emph{Say it:} \emph{kuḷiril nāṉ kuḻantaiyaip pōrvaiyāl pōrttiṉēṉ.} — In the cold I wrapped the child in a blanket.
  \item \emph{Recall:} say the Tamil for an air conditioner, then the Tamil for hot water, then say \emph{pōrttu} again
\end{itemize}

\begin{tcolorbox}[breakable,colback=teal!4,colframe=teal!35!black,title={Before you move on}]
What does \ta{போர்த்து} mean? (\textquotedblleft{}To wrap up; to cover with a cloth\textquotedblright{}.) Say it once more.
\end{tcolorbox}
\section[\texorpdfstring{\ta{இளைப்பாறு} (iḷaippāṟu)}{இளைப்பாறு (iḷaippāṟu)}]{\ta{இளைப்பாறு} (iḷaippāṟu) — to relax; to rest}

\label{lesson:TA-C419-ilaipparu}



\begin{quote}
Before the new one: say the Tamil for hot water, then the Tamil for to wrap up.
\end{quote}

\subsection*{You'll want to know: \ta{இளைப்பாறு}}
\textbf{\ta{இளைப்பாறு}} — \emph{iḷaippāṟu} — \textquotedblleft{}to relax; to rest\textquotedblright{}.

\begin{quote}
\ta{நீண்ட} \ta{பயணத்துக்குப்} \ta{பிறகு} \ta{நாங்கள்} \ta{விடுதியில்} \ta{இளைப்பாறினோம்}.
\end{quote}

\begin{quote}
\emph{nīṇṭa payaṇattukkup piṟaku nāṅkaḷ viṭutiyil iḷaippāṟiṉōm.}
\end{quote}

\begin{quote}
After the long journey we rested at the lodge.
\end{quote}

Past (I …): \ta{இளைப்பாறினேன்}, \emph{iḷaippāṟiṉēṉ}.

\ta{இளை}, to grow tired, + \ta{ஆறு}, to cool down: to recover from tiredness.

\begin{cousinweb}
\textbf{\ta{இளைப்பாறு}} (\emph{iḷaippāṟu}) is built from \ta{இளை}, \textquotedblleft{}to grow weary\textquotedblright{} (DEDR 512), and \ta{ஆறு}, \textquotedblleft{}to cool, to be eased\textquotedblright{} (DEDR 404).
\end{cousinweb}

\begin{grammarlens}[title={What you've built}]
One more word for getting around.
\end{grammarlens}

\subsection*{Your turn}
\begin{itemize}
\raggedright
  \item \emph{Say it:} \emph{iḷaippāṟu}
  \item \emph{Say it:} \emph{iḷaippāṟu}, once more
  \item \emph{Say it:} \emph{nīṇṭa payaṇattukkup piṟaku nāṅkaḷ viṭutiyil iḷaippāṟiṉōm.} — After the long journey we rested at the lodge.
  \item \emph{Recall:} say the Tamil for hot water, then the Tamil for to wrap up, then say \emph{iḷaippāṟu} again
\end{itemize}

\begin{tcolorbox}[breakable,colback=teal!4,colframe=teal!35!black,title={Before you move on}]
What does \ta{இளைப்பாறு} mean? (\textquotedblleft{}To relax; to rest\textquotedblright{}.) Say it once more.
\end{tcolorbox}
\section[\texorpdfstring{\ta{விருந்தோம்பல்} (viruntōmpal)}{விருந்தோம்பல் (viruntōmpal)}]{\ta{விருந்தோம்பல்} (viruntōmpal) — hospitality}

\label{lesson:TA-C419-viruntompal}



\begin{quote}
Before the new one: say the Tamil for to wrap up, then the Tamil for to relax.
\end{quote}

\subsection*{You'll want to know: \ta{விருந்தோம்பல்}}
\textbf{\ta{விருந்தோம்பல்}} — \emph{viruntōmpal} — \textquotedblleft{}hospitality\textquotedblright{}.

\begin{quote}
\ta{அந்தக்} \ta{கிராம} \ta{மக்களின்} \ta{விருந்தோம்பல்} \ta{எங்களுக்குப்} \ta{பிடித்தது}.
\end{quote}

\begin{quote}
\emph{antak kirāma makkaḷiṉ viruntōmpal eṅkaḷukkup piṭittatu.}
\end{quote}

\begin{quote}
We liked the hospitality of those village people.
\end{quote}

\ta{விருந்து}, a guest or a feast, + \ta{ஓம்பல்}, looking after.

\begin{cousinweb}
\textbf{\ta{விருந்தோம்பல்}} (\emph{viruntōmpal}) is built from \ta{விருந்து}, \textquotedblleft{}a guest, a feast\textquotedblright{} (DEDR 5415), and \ta{ஓம்பல்}, \textquotedblleft{}cherishing\textquotedblright{}, from \ta{ஓம்பு} (DEDR 1056).
\end{cousinweb}

\begin{grammarlens}[title={What you've built}]
One more word for getting around.
\end{grammarlens}

\subsection*{Your turn}
\begin{itemize}
\raggedright
  \item \emph{Say it:} \emph{viruntōmpal}
  \item \emph{Say it:} \emph{viruntōmpal}, once more
  \item \emph{Say it:} \emph{antak kirāma makkaḷiṉ viruntōmpal eṅkaḷukkup piṭittatu.} — We liked the hospitality of those village people.
  \item \emph{Recall:} say the Tamil for to wrap up, then the Tamil for to relax, then say \emph{viruntōmpal} again
\end{itemize}

\begin{tcolorbox}[breakable,colback=teal!4,colframe=teal!35!black,title={Before you move on}]
What does \ta{விருந்தோம்பல்} mean? (\textquotedblleft{}Hospitality\textquotedblright{}.) Say it once more.
\end{tcolorbox}
